
\subsubsection{5.1 AIFS Construct Validity (RQ1)}

Table 1 presents the coefficient estimates from the conditional logit model. The AIFS main-effect coefficient is β₁ = \textbf{+0.0246} (z = 5.313, p < 0.001), indicating that responses with higher AIFS scores are preferred in pairwise comparisons. This result holds while controlling for response length (β₂ = +0.0008, z = 8.727, p < 0.001) and semantic cluster fixed effects. The co-significance of both β₁ and β₂ indicates that formatting structure and response length carry independent preference signal.

\textbf{Table 1. Conditional Logit Coefficient Estimates}

\begin{table}[htbp]
\centering
\begin{tabular}{lllllll}
\toprule
Predictor & Coefficient & Std. Error & z & p-value & Odds Ratio & 95\% CI \\
\midrule
ΔAIFS (β₁) & +0.0246 & 0.005 & 5.313 & < 0.001 & 1.0249 & [0.016, 0.034] (coef) \\
Δlength (β₂) & +0.0008 & 9.19e−05 & 8.727 & < 0.001 & 1.0008 & [0.001, 0.001] (coef) \\
ΔAIFS × split (β₄) & −0.0016 & 0.006 & −0.259 & 0.796 & 0.9984 & [0.9861, 1.0108] \\
\bottomrule
\end{tabular}
\end{table}

\textit{Note: 27,034 cluster fixed effects estimated but not shown. Model: ConditionalLogit, Newton optimizer, 80,342 observations, 27,034 groups, mean group size 3.0, max group size 78. Log-likelihood: −33,519.}

\subsubsection{5.2 Main Hypothesis Test (RQ2)}

The interaction coefficient is β₄ = \textbf{−0.0016308} (z = −0.259, p = 0.7958), corresponding to OR = \textbf{0.9984} with 95\% CI [\textbf{0.9861, 1.0108}]. The likelihood ratio test yields LRT statistic = \textbf{0.067} (p = 0.7957). The Wald and LRT tests are consistent.

\textbf{Table 2. Gate Evaluation for β₄}

\begin{table}[htbp]
\centering
\begin{tabular}{llll}
\toprule
Gate & Threshold & Observed & Outcome \\
\midrule
β₄ directional & > 0 & −0.0016 & Not met \\
OR practical effect & ≥ 1.10 & 0.9984 & Not met \\
Wald p-value & < 0.01 & 0.7958 & Not met \\
CI lower bound & > 1.0 & 0.9861 & Not met \\
\bottomrule
\end{tabular}
\end{table}

None of the four pre-registered gate thresholds are met. The point estimate is slightly negative; both Wald and LRT p-values are approximately 0.80. The online-split annotator condition, as operationalized by the helpful-base versus helpful-online split, does not modulate AIFS preference in the direction predicted by the adaptation hypothesis.

\subsubsection{5.3 Precision Analysis (RQ3)}

The 95\% CI for OR is [0.9861, 1.0108], giving a total CI width of \textbf{0.0247} and a half-width of approximately 0.012. This precision derives from the conditional likelihood across 27,034 cluster strata. The upper CI bound of 1.0108 excludes OR ≥ 1.0108 at the 95\% confidence level. The minimum practical effect threshold set a priori (OR = 1.10) lies far outside the CI. These results are consistent with a precision null rather than a power-limited null: the design has sufficient resolution to detect effects of practically meaningful magnitude, and no such effect is present.

Sensitivity analysis across five cosine similarity thresholds (0.75 to 0.90) shows that the OR for β₄ remains within [0.97, 1.03] across all configurations, indicating that the null result is not an artifact of the specific clustering threshold chosen (0.85). A forest plot of β₄ across four model specifications (varying covariates and cluster granularity) similarly shows a stable null.

\subsubsection{5.4 Mechanism Verification Checks}

Five post-estimation checks were conducted to determine whether the null result reflects a model or data pathology.

\begin{table}[htbp]
\centering
\begin{tabular}{lll}
\toprule
Check & Result & Interpretation \\
\midrule
beta4_fitted & Pass & β₄ was estimated; not dropped or degenerate \\
data_variance & Pass & ΔAIFS has non-trivial variance in both splits \\
split_balanced & Pass & Base and online splits are comparably sized \\
clusters_valid & Pass & Clusters contain sufficient within-cluster variation \\
effect_nonzero & Pass & β₁ ≠ 0; AIFS has preference signal \\
\bottomrule
\end{tabular}
\end{table}

All five checks pass. The null on β₄ occurs in a model where β₁ is positive and significant, confirming that the null interaction is not an artifact of a misconfigured model or uninformative features.

\subsubsection{5.5 Optimizer Diagnostics}

BFGS optimization was attempted first and failed to converge on the full 80,342-observation dataset. The failure is attributable to ill-conditioning of the Hessian from 27,034 cluster fixed effects at a median cluster size of 2.8 pairs, which produces a near-singular aggregated Hessian. Newton-Raphson optimization with step-size regularization converged in 14 iterations with a final gradient norm of 3.2 × 10⁻⁸. McFadden R² was not available because the statsmodels ConditionalLogit implementation does not expose a log-likelihood-null attribute (llnull); this does not affect the coefficient estimates.

